\documentclass[aspectratio=169,10pt]{beamer}

\usetheme{Madrid}
\usecolortheme{default}

\usepackage{amsmath}
\usepackage{amssymb}
\usepackage{booktabs}
\usepackage{graphicx}

\graphicspath{{../output/}}

\definecolor{pgoBlue}{RGB}{37,84,150}
\definecolor{pgoGreen}{RGB}{54,130,91}
\definecolor{pgoOrange}{RGB}{191,104,43}

\setbeamercolor{structure}{fg=pgoBlue}
\setbeamercolor{alerted text}{fg=pgoOrange}
\setbeamertemplate{navigation symbols}{}
\setbeamertemplate{caption}[numbered]
\setbeamertemplate{footline}{%
  \leavevmode%
  \hbox{%
    \begin{beamercolorbox}[
      wd=0.32\paperwidth,ht=2.5ex,dp=1.1ex,center
    ]{author in head/foot}%
      \usebeamerfont{author in head/foot}\insertshortauthor
    \end{beamercolorbox}%
    \begin{beamercolorbox}[
      wd=0.48\paperwidth,ht=2.5ex,dp=1.1ex,center
    ]{title in head/foot}%
      \usebeamerfont{title in head/foot}\insertsectionhead
    \end{beamercolorbox}%
    \begin{beamercolorbox}[
      wd=0.20\paperwidth,ht=2.5ex,dp=1.1ex,center
    ]{date in head/foot}%
      \usebeamerfont{date in head/foot}%
      \insertframenumber{} / \inserttotalframenumber
    \end{beamercolorbox}%
  }%
  \vskip0pt%
}

\newcommand{\good}[1]{\textcolor{pgoGreen}{#1}}
\newcommand{\warn}[1]{\textcolor{pgoOrange}{#1}}
\newcommand{\norm}[1]{\left\lVert #1\right\rVert}
\newcommand{\eref}{E_{\mathrm{ref}}}
\newcommand{\Kff}{K_{\!ff}}
\newcommand{\Eup}{E_{up}}
\newcommand{\Eupf}{E_{up,f}}

\title[Equilibrium-Constrained Calibration]{Systematic Poking Calibration through Static Reaction Forces}
\subtitle{Implicit Equilibrium Sensitivity and Positive-Parameter Gauss--Newton}
\author{libpgo material calibration experiment}
\date{July 2026}

\AtBeginSection[]{%
  \begin{frame}[plain]
    \centering
    \vfill
    {\usebeamercolor[fg]{structure}\Large
      Section \insertsectionnumber\ of 7}

    \vspace{1.2em}
    \makebox[\linewidth][c]{%
      \begin{beamercolorbox}[
        wd=0.84\paperwidth,sep=1.5em,center
      ]{title}
        {\LARGE\bfseries\insertsectionhead}
      \end{beamercolorbox}%
    }

    \vspace{1.4em}
    {\large Systematic Poking Calibration through Static Reaction Forces}
    \vfill
  \end{frame}
}

\begin{document}

\begin{frame}
  \titlepage
\end{frame}

\begin{frame}{Roadmap}
  \tableofcontents[hideallsubsections]
\end{frame}

\section{Problem Setup and Notation}

\begin{frame}{Question, Goal, and Scope}
  \begin{block}{Question}
    Can a Systematic Poking material be identified from force--displacement
    observations when every predicted force requires a nonlinear static solve?
  \end{block}

  \vspace{0.35em}
  \begin{columns}[T]
    \begin{column}{0.49\linewidth}
      \textbf{Covered here}
      \begin{itemize}
        \item Nonlinear FEM equilibrium
        \item Scalar reaction-force observations
        \item Implicit differentiation through equilibrium
        \item Positive-parameter Gauss--Newton
        \item Finite-difference and identifiability checks
      \end{itemize}
    \end{column}
    \begin{column}{0.49\linewidth}
      \textbf{Deferred extension}
      \begin{itemize}
        \item Indenter contact
        \item Unknown contact set
        \item Experimental noise
        \item Collapse and inversion
      \end{itemize}
    \end{column}
  \end{columns}
\end{frame}

\begin{frame}{Notation I: Material Parameters}
  \small
  \begin{table}
    \centering
    \begin{tabular}{lll}
      \toprule
      Symbol & Dimension / units & Meaning \\
      \midrule
      $q_k$ & scalar, stress &
      sampled stretch curvature $f''(s_k)$ \\
      $\lambda$ & scalar, stress &
      scale multiplying the fixed volume spline $\widehat h$ \\
      $p$ & $n_p\times1$, stress &
      physical parameter vector $[q_0,\ldots,q_{16},\lambda]^T$ \\
      $n_p$ & scalar integer &
      number of fitted physical parameters; here $n_p=18$ \\
      $\theta$ & $n_p\times1$, dimensionless &
      unconstrained log parameters optimized numerically \\
      $\eref$ & scalar, stress &
      reference modulus used to scale $p$ and force residuals \\
      \bottomrule
    \end{tabular}
  \end{table}

  The conversion between optimizer variables and physical parameters is
  \[
    \boxed{p=\eref\exp(\theta)}
  \]
  with the exponential applied componentwise.
\end{frame}

\begin{frame}{Notation II: FEM State, Energy, and Force}
  \small
  \begin{table}
    \centering
    \begin{tabular}{lll}
      \toprule
      Symbol & Dimension & Meaning \\
      \midrule
      $n_u$ & scalar integer &
      total displacement DOFs; here $n_u=81$ \\
      $u$ & $n_u\times1$ &
      complete nodal displacement vector \\
      $u_f$ & $n_f\times1$ &
      unknown displacement entries on free DOFs \\
      $u_c$ & $n_c\times1$ &
      prescribed displacement entries fixed by one load case \\
      $u_f^\star(p)$ & $n_f\times1$ &
      equilibrium free state obtained by the static solve \\
      $E_h(u,p)$ & scalar &
      assembled discrete hyperelastic energy \\
      $g(u,p)$ & $n_u\times1$ &
      energy gradient $\partial E_h/\partial u$; nodal internal forces \\
      \bottomrule
    \end{tabular}
  \end{table}

  The superscript $\star$ on $u_f^\star$ means ``equilibrium solution,'' not
  target data. The subscript $h$ in $E_h$ denotes the discrete hyperelastic
  energy; it is not Young's modulus $E$.
\end{frame}

\begin{frame}{Notation III: Load Cases, Observations, and Residuals}
  \small
  \begin{table}
    \centering
    \begin{tabular}{lll}
      \toprule
      Symbol & Dimension & Meaning \\
      \midrule
      $j$ & scalar integer &
      load-case index; $j=1,\ldots,N$ within one data split \\
      $c_j$ & $n_u\times1$ &
      selector that sums the measured constrained-force entries \\
      $y_j(p)$ & scalar force &
      predicted reaction after solving equilibrium with parameters $p$ \\
      $y_j^\star$ & scalar force &
      synthetic target reaction for load case $j$ \\
      $r_j(\theta)$ & scalar, normalized &
      fitting error $(y_j(p(\theta))-y_j^\star)/\eref$ \\
      $r(\theta)$ & $N\times1$ &
      residual vector formed by stacking all $r_j$ \\
      \bottomrule
    \end{tabular}
  \end{table}

  \[
    \boxed{
      \theta\longrightarrow p
      \longrightarrow u_f^\star(p)
      \longrightarrow g
      \longrightarrow y_j(p)
      \longrightarrow r_j(\theta)
    }
  \]

  Here $\star$ on $y_j^\star$ means ``target data.'' Later derivations often
  suppress the case index $j$ to reduce clutter.

  Because the reference top-face area is $A_0=1$, the physical force scale
  $\eref A_0$ equals $\eref$ numerically in this unit-cube experiment.
\end{frame}

\begin{frame}{Notation IV: Derivatives and Indexing}
  \small
  Jacobians use \emph{output entries as rows} and \emph{input entries as
  columns}. Let $n_f$ be the number of free DOFs and $n_p=18$.

  \begin{table}
    \centering
    \begin{tabular}{lll}
      \toprule
      Symbol & Dimension & Meaning \\
      \midrule
      $\Kff$ & $n_f\times n_f$ &
      tangent stiffness restricted to free DOFs \\
      $\Eup$ & $n_u\times n_p$ &
      mixed derivative $\partial g/\partial p
      =\partial^2E_h/\partial u\,\partial p$ \\
      $\Eupf$ & $n_f\times n_p$ &
      free rows of $\Eup$ \\
      $du_f^\star/dp$ & $n_f\times n_p$ &
      equilibrium displacement sensitivity \\
      $K_{:f}$ & $n_u\times n_f$ &
      all force rows, free displacement columns \\
      \bottomrule
    \end{tabular}
  \end{table}

  Subscript $f$ or $c$ selects free or constrained rows/columns; a colon
  means ``all.'' Thus $E_{up,f}$ selects free force rows, while $K_{:f}$
  keeps all force rows and selects free displacement columns.
\end{frame}

\begin{frame}{From Direct Stress Fitting to Static Fitting}
  \small
  In the direct constitutive experiment, the deformation is supplied:
  \[
    F_c
    \longrightarrow
    P(F_c;p)
    \longrightarrow
    r_c(p).
  \]
  No state variable needs to be solved.

  \vspace{0.6em}
  Here, only boundary displacement and reaction force are observed:
  \[
    p
    \longrightarrow
    u^\star(p)
    \longrightarrow
    y(p)
    \longrightarrow
    r(p).
  \]

  \begin{alertblock}{The new mathematical difficulty}
    The equilibrium displacement $u^\star$ changes with the material
    parameters.  A correct force Jacobian must include both the direct
    material derivative and this induced state change.
  \end{alertblock}
\end{frame}

\section{Materials and Parameterization}

\begin{frame}{Synthetic Target: Compressible Neo-Hookean}
  \small
  Given Young's modulus $E$ and Poisson ratio $\nu$,
  \[
    \mu=\frac{E}{2(1+\nu)},
    \qquad
    \lambda_{\mathrm{NH}}
    =\frac{E\nu}{(1+\nu)(1-2\nu)}.
  \]

  For principal stretches $s_1,s_2,s_3>0$ and $J=s_1s_2s_3$,
  \[
    \psi_{\mathrm{NH}}
    =
    \sum_{i=1}^{3}
    \left[
      \frac{\mu}{2}(s_i^2-1)-\mu\log s_i
    \right]
    +
    \frac{\lambda_{\mathrm{NH}}}{2}(\log J)^2.
  \]

  \begin{block}{Ground truth and initialization}
    Target: $E^\star=2.0\times10^5,\ \nu^\star=0.35$.
    Initial guess: $E_0=1.2\times10^5,\ \nu_0=0.25$.
  \end{block}
\end{frame}

\begin{frame}{Fitted Model: Systematic Poking}
  \small
  \[
    \psi_{\mathrm{SP}}(s_1,s_2,s_3;p)
    =
    \sum_{i=1}^{3}f(s_i;q)
    +\lambda\,\widehat h(J).
  \]

  The stretch curvature is a piecewise-linear interpolation:
  \[
    f''(s;q)=\sum_{k=0}^{16}q_k\phi_k(s),
    \qquad
    q_k=f''(s_k),
  \]
  integrated with $f(1)=f'(1)=0$.

  The volume shape $\widehat h$ is fixed from log-squared curvature samples:
  \[
    \widehat h''(J_k)=\frac{1-\log J_k}{J_k^2},
    \qquad
    \widehat h(1)=\widehat h'(1)=0.
  \]

  \[
    p=
    \begin{bmatrix}q_0&\cdots&q_{16}&\lambda\end{bmatrix}^{T}
    \in\mathbb R^{18}.
  \]
\end{frame}

\begin{frame}{Positive Physical Parameters}
  \small
  Optimize unconstrained, dimensionless variables $\theta$:
  \[
    p_i=\eref\exp(\theta_i),
    \qquad
    \eref=2.0\times10^5.
  \]
  Therefore,
  \[
    q_k>0,\qquad \lambda>0
  \]
  for every trial step and every line-search point.

  The parameter derivative is diagonal:
  \[
    \frac{\partial p_i}{\partial\theta_j}
    =p_i\delta_{ij},
    \qquad
    \frac{\partial p}{\partial\theta}
    =\operatorname{diag}(p).
  \]

  \begin{block}{Role of $\eref$}
    $\exp(\theta_i)$ is dimensionless. Multiplication by $\eref$ restores
    stress units and keeps typical $\theta_i$ near order one.
  \end{block}
\end{frame}

\section{Load Cases and Observations}

\begin{frame}{Discrete Boundary-Value Problem}
  \small
  The test specimen is a unit cube discretized by a $2\times2\times2$
  cubic-linear mesh:
  \[
    8\ \text{elements},\qquad
    27\ \text{vertices},\qquad
    n_u=81\ \text{displacement DOFs}.
  \]

  Split the displacement vector into free and prescribed components:
  \[
    u=
    \begin{bmatrix}u_f\\u_c\end{bmatrix}.
  \]

  For each load case, $u_c$ is fixed by the boundary condition while $u_f$
  is determined by energy minimization:
  \[
    u_f^\star(p)
    =
    \arg\min_{u_f} E_h(u_f,u_c;p).
  \]

  \begin{alertblock}{Important}
    $u_f^\star$ depends on $p$ even though the prescribed displacement $u_c$
    does not.
  \end{alertblock}
\end{frame}

\begin{frame}{Load Case I: Free Uniaxial}
  \small
  Let $a>0$ be the prescribed axial stretch and let
  $\mathcal T=\{v:Y_v=1\}$ denote the top-face vertices.
  \[
    u_y(Y=0)=0,
    \qquad
    u_y(Y=1)=a-1.
  \]
  Tangential motion is free except for three scalar rigid-mode pins:
  \[
    u_x(v_0)=u_z(v_0)=0,
    \qquad
    u_z(v_1)=0,
    \qquad v_0,v_1\in\{Y=0\}.
  \]
  They remove translation in $x,z$ and rotation about $y$ without fixing
  the bottom face laterally.

  The admissible homogeneous response is
  \[
    F_{\mathrm{free}}=
    \operatorname{diag}(b,a,b),
  \]
  where the static solve selects $b$ from zero lateral traction. For a nearly
  incompressible response, $b\approx a^{-1/2}$ and $J=b^2a\approx1$.

  \begin{block}{Information supplied}
    This case probes axial response together with material-dependent Poisson
    contraction or expansion.
  \end{block}
\end{frame}

\begin{frame}{Free Uniaxial: Why Exactly Three Pins?}
  \small
  After prescribing the top and bottom $y$ displacements, three tangential
  rigid modes remain: $T_x$, $T_z$, and $\omega_y$. Choose bottom vertices
  $v_0,v_1$ with $Z_1=Z_0$ and $X_1\neq X_0$, then impose
  \[
    \boxed{u_x(v_0)=0,\quad u_z(v_0)=0,\quad u_z(v_1)=0.}
  \]
  The first two constraints remove $T_x,T_z$. For a small rotation about
  $y$ around $v_0$,
  \[
    u_x=\omega_y(Z-Z_0),\qquad
    u_z=-\omega_y(X-X_0).
  \]
  Therefore $u_z(v_1)=0$ and $X_1\neq X_0$ imply $\omega_y=0$.

  \begin{block}{Why lateral relaxation is still free}
    A homogeneous transverse stretch $b$ about $v_0$ has
    \[
      u_x=(b-1)(X-X_0),\qquad
      u_z=(b-1)(Z-Z_0).
    \]
    It satisfies all three pins because $Z_1=Z_0$, for any value of $b$.
    Thus the pins select position and orientation, while the static solve
    still determines the Poisson contraction or expansion.
  \end{block}
\end{frame}

\begin{frame}{Fitting Target I: Free-Uniaxial Axial Reaction}
  \small
  Neo-Hookean and Systematic Poking are solved separately under the same
  prescribed stretch $a$:
  \[
    u_{\mathrm{NH,free}}^\star(a)
    =
    \arg\min_{u_f}E_{\mathrm{NH}}(u_f,u_{c,\mathrm{free}}(a)),
  \]
  \[
    u_{\mathrm{SP,free}}^\star(a;p)
    =
    \arg\min_{u_f}E_{\mathrm{SP}}
      (u_f,u_{c,\mathrm{free}}(a);p).
  \]

  The target and prediction are the top-face axial resultants:
  \[
    \boxed{
      y_{\mathrm{free}}^\star(a)
      =
      \sum_{v\in\mathcal T}
      [g_{\mathrm{NH}}(u_{\mathrm{NH,free}}^\star(a))]_{v,y}
    },
  \]
  \[
    \boxed{
      y_{\mathrm{free}}(a;p)
      =
      \sum_{v\in\mathcal T}
      [g_{\mathrm{SP}}(u_{\mathrm{SP,free}}^\star(a;p),p)]_{v,y}
    }.
  \]
  Bottom-face reactions are not included in this scalar observation.
\end{frame}

\begin{frame}{Load Case II: Confined Uniaxial}
  \small
  The axial displacement is prescribed as before:
  \[
    u(x,y=0,z)=0,
    \qquad
    u_y(y=1)=a-1.
  \]
  In addition, lateral displacement is fixed at every vertex:
  \[
    u_x=0,\qquad u_z=0.
  \]
  Thus this is an idealized uniaxial-strain test, stronger than a deformable
  sleeve or penalty contact model. The deformation is approximately
  \[
    \boxed{
      F_{\mathrm{confined}}
      =
      \operatorname{diag}(1,a,1),
      \qquad
      J=a
    }.
  \]

  \begin{block}{Information supplied}
    The material cannot relax laterally, so volume change is forced directly.
    This case strongly excites the volumetric parameter $\lambda$.
  \end{block}
\end{frame}

\begin{frame}{Fitting Target II: Confined Axial Reaction}
  \small
  Separate static solves give
  \[
    u_{\mathrm{NH,conf}}^\star(a),
    \qquad
    u_{\mathrm{SP,conf}}^\star(a;p).
  \]
  The fitted scalar is still only the top-face axial resultant:
  \[
    \boxed{
      y_{\mathrm{conf}}^\star(a)
      =
      \sum_{v\in\mathcal T}
      [g_{\mathrm{NH}}(u_{\mathrm{NH,conf}}^\star(a))]_{v,y}
    },
  \]
  \[
    \boxed{
      y_{\mathrm{conf}}(a;p)
      =
      \sum_{v\in\mathcal T}
      [g_{\mathrm{SP}}(u_{\mathrm{SP,conf}}^\star(a;p),p)]_{v,y}
    }.
  \]

  \begin{alertblock}{Not fitted}
    Enforcing $u_x=u_z=0$ creates lateral constrained reactions, but their
    $x$- and $z$-components are not included in the current loss. Bottom-face
    reactions are also omitted.
  \end{alertblock}
\end{frame}

\begin{frame}{Load Case III: Simple Shear}
  \small
  For shear amount $\gamma$, clamp the bottom face and prescribe
  \[
    u_x(y=1)=\gamma,
    \qquad
    u_y(y=1)=0
  \]
  on the top face. The ideal affine deformation is
  \[
    F_{\mathrm{shear}}(\gamma)
    =
    I+\gamma e_xe_y^T
    =
    \begin{bmatrix}
      1&\gamma&0\\
      0&1&0\\
      0&0&1
    \end{bmatrix},
    \qquad J=1.
  \]

  \begin{block}{Information supplied}
    This case excites a deformation path distinct from normal loading and
    helps separate stretch-curvature parameter directions.
  \end{block}
\end{frame}

\begin{frame}{Fitting Target III: Simple-Shear Tangential Reaction}
  \small
  Interpret the top face as a rigid loading platen with a horizontal load
  cell. The imposed input is its displacement $\gamma$; the measured output
  is the total top-face force in the loading direction $x$.

  For the same $\gamma$, solve two separate equilibrium problems:
  \[
    u_{\mathrm{NH,shear}}^\star(\gamma),
    \qquad
    u_{\mathrm{SP,shear}}^\star(\gamma;p).
  \]
  Their target and predicted measurements are
  \[
    \boxed{
      y_{\mathrm{shear}}^\star(\gamma)
      =
      \sum_{v\in\mathcal T}
      [g_{\mathrm{NH}}(u_{\mathrm{NH,shear}}^\star(\gamma))]_{v,x}
    },
  \]
  \[
    \boxed{
      y_{\mathrm{shear}}(\gamma;p)
      =
      \sum_{v\in\mathcal T}
      [g_{\mathrm{SP}}(u_{\mathrm{SP,shear}}^\star(\gamma;p),p)]_{v,x}
    }.
  \]
  Thus one shear load case contributes the scalar residual
  \[
    \boxed{
      r_{\mathrm{shear}}(\theta;\gamma)
      =
      \frac{
        y_{\mathrm{shear}}(\gamma;p(\theta))
        -y_{\mathrm{shear}}^\star(\gamma)
      }{\eref}
    }.
  \]
\end{frame}

\begin{frame}{Shear Observation: Included and Excluded Forces}
  \small
  \begin{table}
    \centering
    \begin{tabular}{lll}
      \toprule
      Quantity & In loss? & Reason \\
      \midrule
      Top resultant $R_x$ & Yes &
      horizontal load-cell measurement \\
      Bottom resultant $R_x$ & No &
      equal and opposite to top at equilibrium \\
      Top normal resultant $R_y$ & No &
      separate Poynting-effect measurement \\
      Individual nodal reactions & No &
      aggregate force avoids mesh-dependent weighting \\
      \bottomrule
    \end{tabular}
  \end{table}

  Fixing $u_y=0$ on the top may create a nonzero normal reaction $R_y$ even
  though the ideal affine shear has $J=1$. The present experiment deliberately
  fits only the tangential force--displacement curve
  \[
    \gamma\longmapsto R_x(\gamma).
  \]
  Adding $R_y(\gamma)$ would define a second observation and would require an
  explicit relative weight in the loss.

  \begin{block}{Signed shear levels}
    For an isotropic material,
    $R_x(-\gamma)=-R_x(\gamma)$. Using both signs checks material and
    implementation symmetry while balancing the sampled loading directions.
  \end{block}
\end{frame}

\section{Dataset and Objective}

\begin{frame}{Assembling the 36 Training Residuals}
  \small
  \begin{table}
    \centering
    \begin{tabular}{lclc}
      \toprule
      Protocol & Coordinate & Fitted scalar reaction & Train cases \\
      \midrule
      Free uniaxial & $a$ &
      $\displaystyle\sum_{v\in\mathcal T}g_{v,y}$ & 16 \\
      Confined uniaxial & $a$ &
      $\displaystyle\sum_{v\in\mathcal T}g_{v,y}$ & 16 \\
      Simple shear & $\gamma$ &
      $\displaystyle\sum_{v\in\mathcal T}g_{v,x}$ & 4 \\
      \bottomrule
    \end{tabular}
  \end{table}

  For each case $j$,
  \[
    r_j(\theta)
    =
    \frac{y_j(p(\theta))-y_j^\star}{\eref}.
  \]
  The complete training objective is
  \[
    \boxed{
      L(\theta)
      =
      \frac12\sum_{j=1}^{36}r_j(\theta)^2
      +
      \frac{\alpha}{2}\norm{D_2\theta_q}^2
    }.
  \]
\end{frame}

\begin{frame}{What Does One Dataset Record Contain?}
  \small
  One synthetic observation is
  \[
    \boxed{
      \mathcal D_j
      =
      \bigl(
        \text{protocol}_j,\,
        \xi_j,\,
        \text{boundary conditions}_j,\,
        y_j^\star
      \bigr)
    },
  \]
  where $\xi_j=a$ for a uniaxial case and $\xi_j=\gamma$ for shear.
  The target is generated by
  \[
    \text{Neo-Hookean}
    \longrightarrow
    \text{static solve}
    \longrightarrow
    u_j^\star
    \longrightarrow
    y_j^\star.
  \]

  \begin{table}
    \centering
    \begin{tabular}{llrr}
      \toprule
      Split & Protocol & Coordinate & $y_j^\star/\eref$ \\
      \midrule
      Train & Free uniaxial & $a=0.707107$ & $-0.399805$ \\
      Train & Confined uniaxial & $a=0.707107$ & $-0.685460$ \\
      Train & Simple shear & $\gamma=0.4$ & $0.126665$ \\
      Validation & Free uniaxial & $a=0.738413$ & $-0.342510$ \\
      Final test & Simple shear & $\gamma=0.5$ & $0.158874$ \\
      \bottomrule
    \end{tabular}
  \end{table}

  The dataset does not prescribe a full deformation gradient for free
  uniaxial loading; its equilibrium deformation is an output of the solve.
\end{frame}

\begin{frame}{Training Set: Exact Loading Coordinates}
  \small
  The 17 spline knots are
  \[
    s_k=2^{-1+k/8},
    \qquad k=0,\ldots,16,
  \]
  so they are uniform in $\log s$, span $[0.5,2]$, and include $s_8=1$.
  Training uses every non-rest knot:
  \[
    a\in\{s_k:k\neq8\}.
  \]
  Representative values are
  \[
    0.5,\ 0.5453,\ 0.5946,\ldots,\ 0.9170,\ 
    1.0905,\ldots,\ 1.8340,\ 2.0.
  \]

  \begin{table}
    \centering
    \begin{tabular}{lclc}
      \toprule
      Protocol & Coordinates & Observation & Cases \\
      \midrule
      Free uniaxial & $a=s_k,\ k\neq8$ & top $R_y$ & 16 \\
      Confined uniaxial & $a=s_k,\ k\neq8$ & top $R_y$ & 16 \\
      Simple shear & $\gamma=\{-0.8,-0.4,0.4,0.8\}$ &
      top $R_x$ & 4 \\
      \bottomrule
    \end{tabular}
  \end{table}

  Hence $N_{\mathrm{train}}=16+16+4=36$. Only these 36 observations enter
  the optimization of the material parameters $\theta$.
\end{frame}

\begin{frame}{Validation and Final Test: Off-Knot Loading}
  \small
  Validation samples each log-knot interval at its midpoint:
  \[
    a_{k,1/2}=\sqrt{s_ks_{k+1}},
    \qquad k=0,\ldots,15,
  \]
  for both uniaxial protocols, plus
  $\gamma=\{-0.6,-0.2,0.2,0.6\}$. Thus
  \[
    N_{\mathrm{val}}=16\times2+4=36.
  \]

  Final test uses two new positions in every interval:
  \[
    \log a_{k,1/4}
    =\tfrac34\log s_k+\tfrac14\log s_{k+1},
    \qquad
    \log a_{k,3/4}
    =\tfrac14\log s_k+\tfrac34\log s_{k+1},
  \]
  plus
  $\gamma=\{-0.7,-0.5,-0.3,-0.1,0.1,0.3,0.5,0.7\}$, giving
  \[
    N_{\mathrm{test}}=32\times2+8=72.
  \]

  \begin{block}{Different roles}
    Train fits $\theta$; validation selects the smoothness weight $\alpha$;
    final test reports performance after that choice. Off-knot coordinates
    test interpolation rather than memorization at spline knots.
  \end{block}
\end{frame}

\section{Equilibrium Sensitivity}

\begin{frame}{Static Equilibrium}
  \small
  Define the full internal-force vector
  \[
    g(u,p)=\frac{\partial E_h(u,p)}{\partial u}
    \in\mathbb R^{n_u}.
  \]
  Its free-DOF block is
  \[
    g_f(u_f,u_c,p)
    =
    \frac{\partial E_h}{\partial u_f}.
  \]

  The equilibrium state satisfies
  \[
    \boxed{
      g_f(u_f^\star,u_c,p)=0
    }.
  \]

  Forces on free DOFs vanish, but constrained entries of $g$ generally do
  not: those entries are the support reactions required to enforce $u_c$.

  \begin{alertblock}{Notation}
    $g$ is called the energy gradient in the code. Depending on the chosen
    external-force convention, the applied load has the opposite sign.
  \end{alertblock}
\end{frame}

\begin{frame}{Inner Nonlinear Static Solve}
  \small
  At Newton iteration $\ell$, assemble
  \[
    g_f^\ell
    =
    g_f(u_f^\ell,u_c,p),
    \qquad
    \Kff^\ell
    =
    \frac{\partial g_f}{\partial u_f}
    =
    \frac{\partial^2E_h}{\partial u_f^2}.
  \]
  Solve
  \[
    \Kff^\ell\Delta u_f=-g_f^\ell,
    \qquad
    u_f^{\ell+1}=u_f^\ell+\alpha_\ell\Delta u_f,
  \]
  with backtracking.

  The implementation first uses libpgo's \texttt{NewtonOptimizer}, then a
  dense Newton polish until
  \[
    \norm{g_f}_\infty\le 10^{-6}.
  \]

  At $\eref=2\times10^5$, this is a relative force residual below
  $5\times10^{-12}$.
\end{frame}

\begin{frame}{Reaction Observation}
  \small
  For one load case, suppress the index $j$ and write
  \[
    c:=c_j\in\mathbb R^{n_u}.
  \]
  select and sum the constrained force entries associated with the measured
  top-face direction.

  The scalar model prediction is
  \[
    \boxed{
      y(p)=c^T g(u^\star(p),p)
    }.
  \]

  Dimensions are
  \[
    c^T:\ 1\times n_u,\qquad
    g:\ n_u\times1,\qquad
    y:\ 1\times1.
  \]

  \begin{block}{Physical meaning}
    $y$ is the net reaction conjugate to the prescribed top-face
    displacement: axial force for uniaxial loading and horizontal force for
    shear.
  \end{block}
\end{frame}

\begin{frame}{Why the Partial Material Derivative Is Not Enough}
  \small
  Continue to suppress the load-case index $j$.
  The reaction is the composite function
  \[
    y(p)=c^Tg(u^\star(p),p).
  \]
  Therefore,
  \[
    \frac{dy}{dp}
    =
    c^T
    \left(
      \underbrace{\frac{\partial g}{\partial p}}_{\text{direct material effect}}
      +
      \underbrace{\frac{\partial g}{\partial u_f}
      \frac{du_f^\star}{dp}}_{\text{equilibrium state effect}}
    \right).
  \]

  Holding $u^\star$ fixed would retain only the first term.

  \begin{alertblock}{Consequence}
    Two materials can produce different lateral relaxation under the same
    prescribed axial displacement. Ignoring $du_f^\star/dp$ misses the
    resulting change in measured reaction.
  \end{alertblock}
\end{frame}

\begin{frame}{Implicit Differentiation I: Differentiate Equilibrium}
  \small
  Start from
  \[
    g_f(u_f^\star(p),u_c,p)=0.
  \]
  Differentiate both sides with respect to all physical parameters:
  \[
    \frac{\partial g_f}{\partial u_f}
    \frac{du_f^\star}{dp}
    +
    \frac{\partial g_f}{\partial p}
    =0.
  \]

  Introduce
  \[
    \Kff
    :=
    \frac{\partial g_f}{\partial u_f},
    \qquad
    \Eupf
    :=
    \frac{\partial g_f}{\partial p}.
  \]
  Then
  \[
    \boxed{
      \frac{du_f^\star}{dp}
      =
      -\Kff^{-1}\Eupf
    }.
  \]
\end{frame}

\begin{frame}{Stiffness Block Structure: What Is $K_{:f}$?}
  \small
  Order displacement and force entries by free and constrained DOFs:
  \[
    u=
    \begin{bmatrix}u_f\\u_c\end{bmatrix},
    \qquad
    g=
    \begin{bmatrix}g_f\\g_c\end{bmatrix}.
  \]
  The full tangent stiffness is
  \[
    K=\frac{\partial g}{\partial u}
    =
    \begin{bmatrix}
      \dfrac{\partial g_f}{\partial u_f} &
      \dfrac{\partial g_f}{\partial u_c}\\[0.7em]
      \dfrac{\partial g_c}{\partial u_f} &
      \dfrac{\partial g_c}{\partial u_c}
    \end{bmatrix}
    =
    \begin{bmatrix}
      K_{ff}&K_{fc}\\
      K_{cf}&K_{cc}
    \end{bmatrix}.
  \]

  Keeping all force rows but only free-displacement columns gives
  \[
    \boxed{
      K_{:f}
      :=
      \frac{\partial g}{\partial u_f}
      =
      \begin{bmatrix}K_{ff}\\K_{cf}\end{bmatrix}
      \in\mathbb R^{n_u\times n_f}
    }.
  \]
  Thus $K_{:f}\delta u_f$ describes how \emph{all} nodal forces, including
  constrained reactions, change when the free state relaxes.
\end{frame}

\begin{frame}{Reaction Jacobian I: Apply the Chain Rule}
  \small
  Recall
  \[
    y(p)=c^Tg(u^\star(p),p).
  \]
  At fixed prescribed displacement $u_c$,
  \[
    \frac{\partial g}{\partial u_f}=K_{:f},
    \qquad
    \frac{\partial g}{\partial p}=\Eup.
  \]

  Therefore,
  \[
    \frac{dy}{dp}
    =
    c^T
    \left(
      \Eup+
      K_{:f}\frac{du_f^\star}{dp}
    \right).
  \]

  Substitute
  \[
    \frac{du_f^\star}{dp}=-\Kff^{-1}\Eupf.
  \]
\end{frame}

\begin{frame}{Reaction Jacobian II: Final Formula}
  \small
  Substitute the equilibrium sensitivity one line at a time:
  \begin{align*}
    \frac{dy}{dp}
    &=
    c^T\left(
      \Eup+K_{:f}\frac{du_f^\star}{dp}
    \right)\\[0.3em]
    &=
    c^T\left(
      \Eup+K_{:f}\left[-\Kff^{-1}\Eupf\right]
    \right)\\[0.3em]
    &=
    \boxed{
      c^T\left(
        \Eup-K_{:f}\Kff^{-1}\Eupf
      \right)
    }.
  \end{align*}

  \begin{block}{Implementation correspondence}
    \[
      \texttt{free\_sensitivity}
      =-\operatorname{solve}(\Kff,\Eupf),
    \]
    \[
      \texttt{total\_mixed}
      =\Eup+K_{:f}\,\texttt{free\_sensitivity}.
    \]
  \end{block}
\end{frame}

\begin{frame}{Reaction Jacobian III: Constrained-Reaction Form}
  \small
  In this experiment, $c$ selects only constrained reaction entries:
  \[
    c=
    \begin{bmatrix}0\\c_c\end{bmatrix},
    \qquad
    y=c_c^Tg_c.
  \]
  Since
  \[
    c^T\Eup=c_c^TE_{up,c},
    \qquad
    c^TK_{:f}=c_c^TK_{cf},
  \]
  the same derivative becomes
  \[
    \boxed{
      \frac{dy}{dp}
      =
      c_c^T
      \left(
        E_{up,c}
        -
        K_{cf}\Kff^{-1}\Eupf
      \right)
    }.
  \]

  \begin{columns}[T]
    \begin{column}{0.48\linewidth}
      \textbf{Direct material effect}
      \[
        c_c^TE_{up,c}
      \]
      Change material while temporarily holding the state fixed.
    \end{column}
    \begin{column}{0.48\linewidth}
      \textbf{Equilibrium-relaxation effect}
      \[
        -c_c^TK_{cf}\Kff^{-1}\Eupf
      \]
      Restore free equilibrium, which also changes the measured reaction.
    \end{column}
  \end{columns}
\end{frame}

\begin{frame}{Residual and Log-Parameter Jacobian}
  \small
  For load case $j$, normalize the force error by $\eref$:
  \[
    r_j(\theta)
    =
    \frac{y_j(p(\theta))-y_j^\star}{\eref}.
  \]

  Since
  \[
    \frac{\partial p}{\partial\theta}
    =\operatorname{diag}(p),
  \]
  the row of the optimization Jacobian is
  \[
    \boxed{
      \frac{\partial r_j}{\partial\theta}
      =
      \frac{1}{\eref}
      \frac{dy_j}{dp}\operatorname{diag}(p)
    }.
  \]

  This Jacobian contains three derivative layers:
  \[
    \theta
    \longrightarrow p
    \longrightarrow u^\star(p)
    \longrightarrow y(p).
  \]
\end{frame}

\section{Optimization and Verification}

\begin{frame}{Smoothness Regularization}
  \small
  The 17 log-curvatures are
  \[
    \theta_q=
    \begin{bmatrix}\theta_0&\cdots&\theta_{16}\end{bmatrix}^{T}.
  \]
  For uniformly spaced log-stretch knots, let $D_2$ be the second-difference
  operator:
  \[
    (D_2\theta_q)_k
    =
    \theta_{k-1}-2\theta_k+\theta_{k+1}.
  \]

  The objective is
  \[
    L(\theta)
    =
    \frac12\norm{r_{\mathrm{data}}(\theta)}^2
    +
    \frac{\alpha}{2}\norm{D_2\theta_q}^2.
  \]

  Equivalently, append
  \[
    r_{\mathrm{reg}}=\sqrt{\alpha}\,D_2\theta_q,
    \qquad
    J_{\mathrm{reg}}
    =
    \begin{bmatrix}\sqrt{\alpha}\,D_2&0\end{bmatrix}.
  \]
\end{frame}

\begin{frame}{Nested Optimization Pipeline}
  \small
  For every outer evaluation of $\theta$:
  \begin{enumerate}
    \item Convert $\theta$ to positive physical parameters
      $p=\eref\exp(\theta)$.
    \item For every load case, solve
      $g_f(u_f^\star,u_c,p)=0$.
    \item Assemble $y$, $\Kff$, and $\Eup$ at the converged state.
    \item Compute the implicit reaction Jacobian and stack all residual rows.
    \item Append the smoothness residual and Jacobian.
  \end{enumerate}

  The outer damped Gauss--Newton step solves
  \[
    \left[J^TJ+\gamma\,\operatorname{diag}(J^TJ)\right]\Delta\theta
    =
    -J^Tr.
  \]

  A backtracking line search accepts a decrease in the complete objective;
  every candidate step therefore triggers fresh static solves.
\end{frame}

\begin{frame}{Which Optimizer Is Used Where?}
  \small
  There are two different nonlinear problems.

  \begin{columns}[T]
    \begin{column}{0.48\linewidth}
      \textbf{Inner state solve}
      \[
        g_f(u_f,u_c,p)=0.
      \]
      Uses libpgo's
      \texttt{NewtonOptimizer}, fixed damping, and backtracking, followed by
      a dense Newton polish.
    \end{column}
    \begin{column}{0.48\linewidth}
      \textbf{Outer parameter fit}
      \[
        \min_\theta L(\theta).
      \]
      Uses a compact NumPy implementation of damped Gauss--Newton with step
      capping and backtracking.
    \end{column}
  \end{columns}

  \vspace{0.8em}
  \begin{alertblock}{No automatic differentiation through Newton iterations}
    The optimizer differentiates the converged equilibrium equation
    implicitly. It does not backpropagate through the iteration history.
  \end{alertblock}
\end{frame}

\begin{frame}{Full-Pipeline Finite-Difference Check}
  \small
  For parameter direction $e_i$, central finite differences use
  \[
    J_i^{\mathrm{FD}}
    =
    \frac{
      r(\theta+\varepsilon e_i)
      -
      r(\theta-\varepsilon e_i)
    }{2\varepsilon}.
  \]

  Each $r(\theta\pm\varepsilon e_i)$ is evaluated by:
  \[
    \text{new parameters}
    \rightarrow
    \text{fresh static solves}
    \rightarrow
    \text{new reactions}.
  \]

  A six-parameter diagnostic problem gives
  \[
    \frac{\norm{J^{\mathrm{analytic}}-J^{\mathrm{FD}}}_F}
         {\norm{J^{\mathrm{FD}}}_F}
    =
    3.17\times10^{-10},
  \]
  with maximum column-relative error $6.72\times10^{-10}$.

  \begin{block}{What this verifies}
    Constitutive parameter derivatives, FEM assembly, equilibrium
    sensitivity, reaction extraction, and the log-parameter chain rule.
  \end{block}
\end{frame}

\begin{frame}{Identifiability at the Fitted Solution}
  \small
  Let $J_{\mathrm{data}}\in\mathbb R^{36\times18}$ be the training reaction
  Jacobian before regularization.

  For the default 17-knot experiment,
  \[
    \operatorname{rank}(J_{\mathrm{data}})=18,
    \qquad
    \kappa_2(J_{\mathrm{data}})\approx458.
  \]

  Its largest and smallest singular values are
  \[
    \sigma_{\max}=2.048,
    \qquad
    \sigma_{\min}=4.47\times10^{-3}.
  \]

  \begin{block}{Interpretation}
    All 18 local parameter directions are excited by the selected protocols,
    but sensitivity is anisotropic. Smoothness regularization stabilizes the
    weak directions rather than repairing an exact rank deficiency.
  \end{block}
\end{frame}

\begin{frame}{Validation Selection of Smoothness Weight}
  \small
  \begin{table}
    \centering
    \begin{tabular}{rrrr}
      \toprule
      $\alpha$ & Train RMSE & Validation RMSE &
      $\norm{D_2\theta_q}^2$ \\
      \midrule
      $0$ & $2.25{\times}10^{-4}$ & $4.01{\times}10^{-4}$ &
      $7.24{\times}10^{-2}$ \\
      $10^{-4}$ & $2.36{\times}10^{-4}$ & $2.62{\times}10^{-4}$ &
      $1.38{\times}10^{-3}$ \\
      $3{\times}10^{-4}$ & $2.41{\times}10^{-4}$ &
      $2.54{\times}10^{-4}$ & $9.06{\times}10^{-4}$ \\
      $\mathbf{10^{-3}}$ &
      $\mathbf{2.47{\times}10^{-4}}$ &
      $\mathbf{2.49{\times}10^{-4}}$ &
      $\mathbf{6.98{\times}10^{-4}}$ \\
      $3{\times}10^{-3}$ & $2.54{\times}10^{-4}$ &
      $2.50{\times}10^{-4}$ & $6.26{\times}10^{-4}$ \\
      $10^{-2}$ & $2.66{\times}10^{-4}$ & $2.57{\times}10^{-4}$ &
      $5.85{\times}10^{-4}$ \\
      \bottomrule
    \end{tabular}
  \end{table}

  Validation selects $\alpha=10^{-3}$. Final-test error is not part of this
  selection rule.
\end{frame}

\section{Results and Outlook}

\begin{frame}{Quantitative Fit}
  \small
  The selected run terminates by gradient tolerance after seven outer
  Gauss--Newton iterations.

  \begin{table}
    \centering
    \begin{tabular}{lrr}
      \toprule
      Split & Initial normalized RMSE & Fitted normalized RMSE \\
      \midrule
      Training & $3.41\times10^{-1}$ & $2.47\times10^{-4}$ \\
      Validation & $3.05\times10^{-1}$ & $2.49\times10^{-4}$ \\
      Sealed holdout & $3.06\times10^{-1}$ & $2.35\times10^{-4}$ \\
      \bottomrule
    \end{tabular}
  \end{table}

  \[
    L_{\mathrm{final}}=1.45\times10^{-6},
    \qquad
    \min_i p_i=9.22\times10^4>0.
  \]

  The maximum final free-equilibrium residual over all splits is below
  $9.63\times10^{-7}$.
\end{frame}

\begin{frame}{Diagnostics}
  \centering
  \includegraphics[width=0.95\linewidth,height=0.78\textheight,keepaspectratio]
    {fit_diagnostics.png}

  \vspace{-0.2em}
  \small
  Reaction curves agree on train, validation, and final-test cases;
  the loss decreases monotonically under the accepted line-search steps.
\end{frame}

\begin{frame}{Resolution Sweep}
  \small
  A separate unregularized sweep gives:
  \begin{table}
    \centering
    \begin{tabular}{rrrrr}
      \toprule
      Knots & Parameters & Rank & Condition & Validation RMSE \\
      \midrule
      5 & 6 & 6 & 45.4 & $7.45\times10^{-3}$ \\
      9 & 10 & 10 & 123.2 & $1.28\times10^{-3}$ \\
      17 & 18 & 18 & 460.5 & $4.01\times10^{-4}$ \\
      \bottomrule
    \end{tabular}
  \end{table}

  More knots reduce approximation error in this experiment, while the
  Jacobian becomes less well conditioned.

  \begin{alertblock}{Interpretation limit}
    Knot-aligned training observations increase with resolution. This is a
    joint capacity-and-excitation check, not a controlled fixed-data
    capacity ablation.
  \end{alertblock}
\end{frame}

\begin{frame}{Validity Audit}
  \small
  \begin{columns}[T]
    \begin{column}{0.49\linewidth}
      \textbf{\good{Supported}}
      \begin{itemize}
        \item Full analytic Jacobian matches fresh-solve finite differences.
        \item Static residuals are tightly controlled.
        \item The selected data Jacobian is full column rank.
        \item Prediction transfers to off-knot final-test cases.
        \item Positivity is maintained throughout optimization.
      \end{itemize}
    \end{column}
    \begin{column}{0.49\linewidth}
      \textbf{\warn{Not established}}
      \begin{itemize}
        \item Robustness to experimental noise
        \item Mesh-discretization mismatch
        \item Statistical variability across targets
        \item Identifiability under contact
        \item Generalization outside the sampled strain range
      \end{itemize}
    \end{column}
  \end{columns}
\end{frame}

\begin{frame}{Experimental Caveats}
  \small
  \begin{itemize}
    \item The experiment is deterministic and synthetic; there are no
      repeated noisy trials or confidence intervals.
    \item Target and fitted reactions use the same mesh and FEM
      discretization. The test isolates model representation and inverse
      implementation, not discretization robustness.
    \item The regularization sweep is exploratory model selection. The
      holdout is protected from direct selection, but only one target
      material is tested.
    \item A dense Newton polish was added because the backend's relative
      stopping rule alone did not guarantee sufficiently small absolute
      residuals for sensitivity verification.
  \end{itemize}

  \begin{block}{Current confidence}
    High confidence in the implemented equilibrium derivative and synthetic
    calibration pipeline; limited evidence for real poking data.
  \end{block}
\end{frame}

\begin{frame}{Next Highest-Information Experiment}
  \small
  Replace the prescribed top face with a cylindrical indenter and measure a
  force--depth curve:
  \[
    p
    \longrightarrow
    \text{contact equilibrium}
    \longrightarrow
    y_{\mathrm{indenter}}(p).
  \]

  Retain the same verification hierarchy:
  \begin{enumerate}
    \item Tight equilibrium residual for every load step.
    \item Analytic-versus-finite-difference reaction Jacobian.
    \item Singular values and rank of the data Jacobian.
    \item Train/validation/holdout indentation depths.
    \item Mesh-resolution and noise sensitivity.
  \end{enumerate}

  \begin{alertblock}{New risk}
    Contact activation changes the governing branch. Derivatives must be
    checked away from, and then carefully near, contact-set transitions.
  \end{alertblock}
\end{frame}

\begin{frame}{Conclusions}
  \small
  \begin{enumerate}
    \item Static fitting introduces an implicit state map
      $p\mapsto u^\star(p)$.
    \item Differentiating free equilibrium gives
      \[
        \frac{du_f^\star}{dp}=-\Kff^{-1}\Eupf.
      \]
    \item The complete scalar reaction derivative is
      \[
        \frac{dy}{dp}
        =
        c^T\left(\Eup-K_{:f}\Kff^{-1}\Eupf\right).
      \]
    \item The full implementation matches finite differences to about
      $5\times10^{-10}$ relative error.
    \item The 18-parameter fit reduces synthetic final-test RMSE from
      $3.06\times10^{-1}$ to $2.35\times10^{-4}$.
  \end{enumerate}

  \vspace{0.3em}
  \begin{block}{Takeaway}
    The equilibrium-constrained calibration layer is operational and
    verified. Contact-based poking is now the natural next extension.
  \end{block}
\end{frame}

\end{document}
